\subsection{FUNDAMENTACIÓN}

En la actualidad, la biometría de voz se ha consolidado como un mecanismo de autenticación robusto,
permitiendo el control de acceso sin contacto y mejorando la seguridad en múltiples entornos.
Tradicionalmente, estos sistemas han dependido de arquitecturas basadas en la nube o de procesadores
de alto rendimiento debido a la complejidad computacional que requiere la extracción de características
acústicas y la inferencia mediante modelos de aprendizaje automático. Sin embargo, el auge del paradigma
Edge Computing y, específicamente, del TinyML (Machine Learning en el borde), ha posibilitado la
migración de estos algoritmos hacia dispositivos con recursos más limitados (\cite{rayen}). 
En este marco, la combinación de microcontroladores de bajo costo y consumo, junto con micrófonos 
con tecnología MEMS (Sistemas Microelectromecánicos), representa una solución de hardware óptima y 
escalable para la implementación de sistemas de bioacceso autónomos e integrables en el ecosistema 
de la Internet de las Cosas (IoT) (\cite{warden}).

Si bien la literatura documenta avances significativos en el reconocimiento de locutor (SR) 
(\cite{fenu}) , existe un vacío crítico en la intersección de la optimización de hardware
y la sociolingüística aplicada. La inmensa mayoría de los conjuntos de datos vocales (corpus)
y los modelos pre entrenados disponibles comercial o académicamente están sesgados hacia idiomas
hegemónicos, como el inglés, o hacia variantes estandarizadas del español, como  español 
latinoamericano "neutro". La variante del castellano rioplatense posee características
fonológicas, prosódicas y de articulación únicas, tales como el yeísmo y la aspiración de la /s/ 
implosiva (\cite{colantoni}). Al cuantizar (conversión de 32 bits a 8 bits) y comprimir redes 
neuronales para que operen dentro de las limitaciones de memoria (SRAM/Flash) y procesamiento 
de un microcontrolador, los modelos tienden a perder precisión, y esta degradación afecta de 
manera desproporcionada a las variantes dialectales subrepresentadas en el entrenamiento original. 
Actualmente, no existe un marco de desarrollo estandarizado que garantice la eficacia de sistemas 
de bioacceso offline sobre hardware restrictivo calibrado específicamente para la fonética rioplatense.

A partir de la problemática expuesta, la presente investigación se articula en torno a la 
siguiente pregunta: ¿De qué manera se puede ajustar y optimizar un modelo de 
aprendizaje automático para la detección y verificación del hablante, implementado 
en microcontroladores de bajos recursos mediante micrófonos MEMS, que garantice un alto 
nivel de precisión y seguridad operando exclusivamente con locutores de dialecto castellano rioplatense?

El desarrollo de esta investigación se justifica por su doble aporte: tecnológico y 
sociocultural. Desde la perspectiva tecnológica, el proyecto llenará un punto vacío 
en la ingeniería de sistemas embebidos al proponer una metodología documentada para 
la extracción de características acústicas (como coeficientes MFCC), el aumento y 
tratamiento de datos para mitigar el ruido inherente a los micrófonos MEMS (\cite{ko2015}), 
y el diseño de arquitecturas de inferencia ligeras que logren alta precisión frente a variaciones 
fonéticas locales, superando las limitaciones intrínsecas del hardware. Desde una 
perspectiva sociocultural y de usabilidad, este trabajo combate el sesgo lingüístico 
en la Inteligencia Artificial. Al desarrollar tecnología de acceso biométrico que 
reconozca y valide las particularidades del castellano rioplatense, se democratiza 
el acceso a la seguridad tecnológica de vanguardia, garantizando que los usuarios 
locales interactúen con dispositivos que se adaptan a su identidad lingüística, 
y no a la inversa.

% \subsection{INTRODUCCIÓN}
% Históricamente, la complejidad computacional requerida para la extracción de
% características acústicas y la inferencia de modelos obligaba a que estos sistemas dependieran
% de la nube o de procesadores de alto rendimiento. La integración de microcontroladores de
% bajo costo junto con micrófonos MEMS representa hoy una solución escalable y eficiente
% para implementar sistemas de reconocimiento de voz autónomos, preservando la privacidad al
% procesar los datos de forma local en el ecosistema IoT (Barovic \& Moin, 2025).

% Para dar respuesta a esta problemática, el presente trabajo propone diseñar, entrenar y
% optimizar un sistema de control de acceso biométrico por voz en hardware de bajo costo,
% orientado exclusivamente a locutores de dialecto castellano rioplatense. La investigación
% plantea un diseño que procesará la información de forma estrictamente local utilizando un
% microcontrolador de la familia ESP32 y un micrófono digital MEMS. Metodológicamente, se
% abarcará desde la extracción de características acústicas basadas en coeficientes MFCC hasta
% el entrenamiento, la cuantización y la implementación del modelo de red neuronal en el
% dispositivo.

% Finalmente, el desempeño y la viabilidad del sistema se validarán a través de métricas
% específicas como la Tasa de Igual Error (EER), comparando el rendimiento del modelo
% original frente a su versión cuantizada embebida. De este modo, la investigación no solo
% ofrece un aporte tecnológico sustancial para la ingeniería de sistemas embebidos, sino que
% realiza una contribución sociocultural al mitigar el sesgo lingüístico en la IA, permitiendo que
% los dispositivos de seguridad se adapten a la identidad de los usuarios locales, y no a la
% inversa

\subsection{OBJETIVOS}

\subsubsection{Objetivo general}
En lo que compete a esta investigación, el objetivo general de la misma abarca desarrollar
un sistema destinado al control de acceso de personas por medio de reconocimiento de voz en
castellano rioplatense, implementado en hardware de bajo costo

\subsubsection{Objetivos específicos}
\begin{itemize}
    \item  Implementar un algoritmo para la extracción del descriptor MFCC (Mel-Frecuency
Cespectral Coefficient).
    \item Ajuste fino (FT) en PC para su posterior cuantización.
    \item  Evaluar la tasa de error (EER) post cuantización.
    \item Transferir el modelo entrenado al microcontrolador.
    \item Comparar el EER obtenido entre el dispositivo y la PC.
    \item Confeccionar una interfaz de usuario para que el administrador del sistema pueda
generar nuevos enrolamientos de usuarios.
    \item Realizar pruebas de concepto con usuarios que no se hayan empleado para el
entrenamiento del modelo.
\end{itemize}

\subsection{ESTRUCTURA DE LA INVESTIGACIÓN}
\textcolor{red}{En este sección ira detallado como se dividen los distintos apartados de la investigación.}